% WORKING TEMPLATE ONLY. Replace every DRAFT / STUDENT field before submission.
% Compile with the official IEEEtran class (e.g. Overleaf IEEE Conference template).
% Complete the method and evaluation sections before final submission.
\documentclass[conference]{IEEEtran}
\usepackage[T1]{fontenc}
\usepackage[utf8]{inputenc}
\usepackage{amsmath,amssymb}
\usepackage{graphicx}
\usepackage{booktabs}
\usepackage{array}
\usepackage{xcolor}
\usepackage{url}
\newcommand{\draftnote}[1]{\par\noindent\textbf{To be completed:} #1\par}

\title{SudokuScan: Photograph-to-Grid Recognition}
\author{\IEEEauthorblockN{STUDENT NAME}
\IEEEauthorblockA{Department of Computer Engineering, MEF University\\
COMP 482, Fall 2026\\Student number: STUDENT NUMBER}}

\begin{document}
\maketitle
\begin{abstract}
\draftnote{Write the problem, your own outer-corner method, dataset size, and
measured results after completing the experiments. The supplied synthetic smoke
test is not evidence of personal-dataset or hidden-test performance.}
\end{abstract}

\section{Introduction}
\draftnote{Describe the photographed printed-grid problem, its geometric and
recognition challenges, your contribution, and your measured results.}

\section{Corner Detection}
\subsection{Preprocessing and Harris Candidates}
The support implementation resizes large images before detection, compensates
for a smooth illumination field, and evaluates \texttt{cv2.cornerHarris}.
Relative response thresholding and non-maximum suppression produce discrete candidates.
\draftnote{Describe the preprocessing you actually retained. State the working
resolution, size-dependent kernels, actual blockSize, ksize, k, threshold, and NMS radius.}

\subsection{Outer-Corner Selection}
\draftnote{Explain the automatic rule, why it
rejects clutter and page boundaries, how it orders TL/TR/BR/BL, and when it fails.
The current selection function uses contour region proposals, Harris candidate indices and a 10-by-10 line score. Verify the method and record its actual authorship.}

\subsection{Parameter Study}
\draftnote{Use at least six personal images. Insert candidate counts and corner
errors while varying blockSize, ksize, k and response threshold. Explain the values
you chose using your own measurements. Do not replace failed detections with GT.}

\section{Rectification}
An affine map has six parameters,
\begin{equation}
\begin{bmatrix}u\\v\end{bmatrix}=A
\begin{bmatrix}x\\y\\1\end{bmatrix},\quad A\in\mathbb{R}^{2\times3}.
\end{equation}
Three non-collinear point pairs supply six scalar equations. A3 uses TL, TR and BL.
A4 fits all four corners by least squares. A projective map has the form
\begin{equation}
\lambda\begin{bmatrix}u\\v\\1\end{bmatrix}=H
\begin{bmatrix}x\\y\\1\end{bmatrix},\quad H\in\mathbb{R}^{3\times3}.
\end{equation}
Its common scale is unobservable, leaving eight degrees of freedom; four pairs
in general position determine it. The implementation uses the continuous target
boundary $(0,0),(450,0),(450,450),(0,450)$ and a $450\times450$ raster.
\begin{equation}
e_{A3}=\|A_3p_{BR}-(450,450)^T\|_2,
\quad e_{A4}=\sqrt{\frac14\sum_{i=1}^{4}\|A_4p_i-q_i\|_2^2}.
\end{equation}
\draftnote{Verify these statements against your final code, add your residual
measurements and H4 numerical check, and explain parallelism/vanishing points.
Document the explicit four-point LS refit following estimateAffine2D.}

\section{Affine Versus Homography}
\draftnote{Complete measurements on at least eight personal photographs: define obliqueness;
measure A3 fourth-corner/A4 RMS residual; locate all ten lines per direction;
measure line-position, spacing and straightness errors; run identical OCR on
A3/A4/H4; plot the results and interpret the evidence. Explain your conclusion
using the course's weak-perspective model. Include the three rectifications of
your most oblique actual photograph side by side. No conclusions are pre-filled.}

\section{Recognition}
\draftnote{Describe your final margin and empty/filled criterion and justify them
on real photographs. The current implementation uses an 8\% cell margin, connected-component
area/height thresholds, edge-remnant rejection, grayscale-preserving digit
isolation, centering and white padding. Treat defaults as hypotheses, not validated results.}
Grayscale closing estimates the full-grid illumination field before cell thresholding.
The OCR support uses Tesseract through pytesseract. Each filled cell is a separate
page in a multi-page TIFF. The default segmentation mode is PSM 10 and the
character whitelist is 1--9. Empty cells are excluded before OCR. An unread filled
cell receives one twofold retry; raw and prepared experiment arms use the same rule.
\draftnote{Report raw/prepared accuracy, PSM 7/8/10/13 measurements, empty-cell F1,
and common confusions with actual crops. Record engine version and runtime.}

\section{Results}
\draftnote{Fill both development and personal-set results, and C1--C5 breakdowns.
Account for every image, including failures. State the denominator for corner
error when corners are missing. Report empty-positive and filled-positive F1 if needed.}
\begin{table}[ht]
\centering
\caption{Actual evaluation results (not yet measured).}
\begin{tabular}{lcc}
\toprule
Metric & Development & Personal\\
\midrule
Corner error (\% side) & -- & --\\
Cell accuracy & -- & --\\
Empty-cell F1 & -- & --\\
Grid exact match & -- & --\\
Robustness / 20 s & -- & --\\
\bottomrule
\end{tabular}
\end{table}

\section{Failure Analysis}
\draftnote{Include at least four real failure cases with image IDs, condition,
debug crop, numerical error, and a diagnosis of the failing stage. Do not invent
failures or present synthetic tests as phone photographs.}

\section{Conclusion and Limitations}
\draftnote{Write after completing your experiments. Summarize your own findings,
real limitations and remaining work. Check the six-page main-report limit.}

\begin{thebibliography}{4}
\bibitem{opencvfeatures} OpenCV, ``Feature Detection.'' [Online]. Available:
\url{https://docs.opencv.org/4.x/dd/d1a/group__imgproc__feature.html}
\bibitem{opencvwarp} OpenCV, ``Geometric Image Transformations.'' [Online]. Available:
\url{https://docs.opencv.org/4.x/da/d54/group__imgproc__transform.html}
\bibitem{opencvcalib} OpenCV, ``Camera Calibration and 3D Reconstruction.'' [Online]. Available:
\url{https://docs.opencv.org/4.x/d9/d0c/group__calib3d.html}
\bibitem{tesseract} Tesseract, ``Command Line Usage.'' [Online]. Available:
\url{https://tesseract-ocr.github.io/tessdoc/Command-Line-Usage.html}
\end{thebibliography}

\appendices
\section{Dataset and Additional Measurements}
\draftnote{List the original image IDs, puzzle/source IDs, C1--C5 conditions,
capture dates and any additional measurements. Keep phone files unedited with EXIF.}

\section{Reproducibility}
\draftnote{Record final CPU hardware, operating system, dependency versions,
commands, runtime distribution and code commit. Include additional experiment
settings if needed; do not invent a development history.}

\section{AI-Use Disclosure}
\draftnote{This must be the last appendix. Verify and expand the table using
docs/AI\_KULLANIM\_BEYANI.md. Clarify the permitted scope of AI-generated implementation
code with the instructor. Never write `None' for this package.}
\begin{table*}[ht]
\centering
\small
\begin{tabular}{p{0.10\textwidth}p{0.20\textwidth}p{0.28\textwidth}p{0.30\textwidth}}
\toprule
Tool & Purpose / stage & What was requested & What was kept, changed or rejected\\
\midrule
OpenAI Codex & Starter implementation and tools & Preprocessing, Harris/NMS support,
transforms, OCR, validity, interfaces, tests and documentation & AI-generated starter;
student must record actual use and modifications and verify permitted scope.\\
OpenAI Codex & Stage 2 & Outer-corner selection and tests &
AI-generated implementation retained; this does not meet the independent-authorship condition in the brief.\\
OpenAI Codex & Stage 4 & Cross-model comparison analysis &
Not generated; the analysis remains incomplete.\\
\bottomrule
\end{tabular}
\end{table*}
\draftnote{After it is true, personally complete the confirmation in Appendix C
of the assignment: authorship of corner-selection logic and comparison analysis,
and ability to explain and modify every part. No confirmation or signature is
provided by this template.}
Name: \underline{\hspace{25mm}}\quad Student no.: \underline{\hspace{20mm}}\\
Signature: \underline{\hspace{22mm}}\quad Date: \underline{\hspace{20mm}}
\end{document}
